\documentclass[12pt, ukrainian]{article}
\usepackage[a4paper]{geometry}
\setlength\parindent{0mm}
\geometry{top=1.1cm,bottom=1.1cm,left=1.0cm,right=1.0cm}
\pagestyle{empty}
\usepackage[T2A]{fontenc}
\usepackage[utf8]{inputenc}
\usepackage[ukrainian]{babel}
\usepackage{graphicx}
\usepackage{caption}
\DeclareCaptionFormat{nocaption}{}
\captionsetup{format=nocaption,aboveskip=0pt,belowskip=0pt}
\usepackage[Export]{adjustbox}
\adjustboxset{max size={0.9\linewidth}{0.9\paperheight}}
\begin{document}
\begin {Huge}
{{04.12.2023 Контрольна робота \No 2-bis, варіант  \No \textbf { 11 }\par}}

\begin{center}Частина 2
\end{center}
\vspace{10mm}
\end {Huge}

{
\begin {Large}
Якість коду та економність обчислень оцінюється по обом завданням.

\vspace{5mm}

{\begin {center}\textbf{Завдання 1}\end {center}}\par
Написати функцію, що для інтервалу $[a, b)$ знаходить найближчу пару чисел з цього інтервалу, що мають парну кількість різних простих дільників. Якщо таких пар кілька, то цікавить пара з найменшим добутком.

 Використання рядків та структур даних заборонено.



\vspace{10mm}\par
{\begin {center}\textbf{Завдання 2}\end {center}}\par

Написати функцію, що виводить послідовність чисел
$$\scalebox{1.2}{$(-1)^{k}C_{3k-1}^{k}$}, \quad k=2,\ldots,n$$
Оцінюється економність обчислень. Використовувати **, вкладені цикли та виклики функцій \textbf{заборонено}.\par
\end{Large}\newpage

%end



\end{document}

